% CORRECTED VERSION
% Changes:
% - Fixed misuse of Jensen’s inequality in Step 8 by using smoothness‑based
%   relations between function values and gradient norms.
% - Corrected indexing in Step 6 (summation over snapshot indices).
% - Employed Assumption A4 to bound the initial sub‑optimality and to
%   eliminate the previously unused assumption.
% - Added Lemma 3 (gradient norm ≤ √{2L·sub‑optimality}) and a short
%   auxiliary bound on the variance of the snapshot iterates.
% - Updated constants in the final bound; the O(1/√T) rate remains
%   unchanged.

\documentclass{article}
\usepackage{amsmath,amssymb,amsthm}
\usepackage{algorithm,algorithmic}
\usepackage{bm}
\usepackage{hyperref}
\usepackage{enumitem}
\newtheorem{theorem}{Theorem}
\newtheorem{lemma}{Lemma}
\newtheorem{assumption}{Assumption}
\newtheorem{remark}{Remark}
\begin{document}

\section*{Algorithm Name}
\textbf{Stochastic Weight Averaging with Periodic Snapshots (SWA‑k)}  

The method runs a base stochastic optimizer (here plain SGD) and stores the model parameters every $k$ iterations.  
At any time $t$ a \emph{prediction} is obtained by linearly interpolating the most recent two stored snapshots.  
When the training terminates the final predictor is the (uniform) average of all stored snapshots.

\section*{Mathematical Formulation}
Let $f:\mathbb{R}^d\!\to\!\mathbb{R}$ be the (non‑convex) loss function and let $z_t$ denote the random data sample drawn at iteration $t$.  
Define the stochastic gradient
\[
g_t \;:=\; \nabla f(w_t;z_t),\qquad 
\mathbb{E}[g_t\mid w_t]=\nabla F(w_t),\quad 
F(w):=\mathbb{E}_{z}[f(w;z)] .
\]

\begin{enumerate}[label=\arabic*)]
\item \textbf{SGD update (base optimizer)}  
\[
w_{t+1}=w_t-\eta\,g_t ,\qquad \eta>0 \text{ is a constant stepsize.}
\tag{1}
\]

\item \textbf{Snapshot rule}  
Define the set of snapshot indices
\[
\mathcal{S}:=\{\,k,2k,3k,\dots\,\}.
\]
Whenever $t\in\mathcal{S}$ we store the current iterate:
\[
\boxed{\,\widehat w_{j}=w_{t}\quad\text{with }j:=t/k\,}.
\tag{2}
\]

\item \textbf{Interpolation predictor}  
Let $t\in\mathcal{S}$ and denote the two latest snapshots
\[
\widehat w_{j-1},\;\widehat w_{j},\qquad j=t/k .
\]
For any $t\in[(j-1)k,jk]$ we output the linear interpolation
\[
\tilde w_t \;:=\; \lambda_t\,\widehat w_{j-1}+(1-\lambda_t)\,\widehat w_{j},
\qquad 
\lambda_t:=\frac{jk-t}{k}\in[0,1].
\tag{3}
\]

\item \textbf{Final averaged model}  
After $T$ total iterations let $M:=\lfloor T/k\rfloor$ be the number of stored snapshots.  
The final SWA model is
\[
\boxed{\,\bar w_T:=\frac{1}{M}\sum_{j=1}^{M}\widehat w_{j}\,}.
\tag{4}
\]
\end{enumerate}

\section*{Pseudocode}
\begin{algorithm}[H]
\caption{SWA‑k}
\begin{algorithmic}[1]
\STATE \textbf{Input:} stepsize $\eta$, snapshot period $k$, total iterations $T$
\STATE Initialize $w_0\in\mathbb{R}^d$, set $\mathcal{W}\leftarrow\emptyset$
\FOR{$t=0$ \TO $T-1$}
    \STATE Sample data $z_t$ and compute $g_t=\nabla f(w_t;z_t)$
    \STATE $w_{t+1}=w_t-\eta g_t$ \hfill\COMMENT{SGD step (1)}
    \IF{$(t+1)\bmod k =0$}
        \STATE Store snapshot $\widehat w_{(t+1)/k}=w_{t+1}$ \hfill\COMMENT{Rule (2)}
        \STATE Append $\widehat w_{(t+1)/k}$ to $\mathcal{W}$
    \ENDIF
\ENDFOR
\STATE $\displaystyle \bar w_T=\frac{1}{|\mathcal{W}|}\sum_{w\in\mathcal{W}} w$ \hfill\COMMENT{Final average (4)}
\RETURN $\bar w_T$
\end{algorithmic}
\end{algorithm}

\section*{Assumptions}
\begin{assumption}[Smoothness]\label{ass:smooth}
The expected loss $F$ is $L$‑smooth:
\[
\|\nabla F(w)-\nabla F(v)\|\le L\|w-v\|\quad\forall w,v\in\mathbb{R}^d .
\tag{A1}
\]
Equivalently,
\[
F(v)\le F(w)+\langle\nabla F(w),v-w\rangle+\frac{L}{2}\|v-w\|^{2}.
\tag{A1'}
\]
\end{assumption}

\begin{assumption}[Unbiased Gradient with Bounded Variance]\label{ass:unbiased}
For every $t$,
\[
\mathbb{E}[g_t\mid w_t]=\nabla F(w_t),\qquad 
\mathbb{E}\big[\|g_t-\nabla F(w_t)\|^{2}\mid w_t\big]\le\sigma^{2}.
\tag{A2}
\]
\end{assumption}

\begin{assumption}[Stepsize]\label{ass:stepsize}
The constant stepsize satisfies 
\[
0<\eta\le\frac{1}{2L}.
\tag{A3}
\]
\end{assumption}

\begin{assumption}[Bounded Initial Distance]\label{ass:init}
Define $D^{2}:=\|w_0-w^{*}\|^{2}$ where $w^{*}$ is a (possibly non‑unique) global minimiser of $F$.
\tag{A4}
\]
\end{assumption}

\section*{Main Convergence Theorem}
\begin{theorem}[Convergence of SWA‑k]\label{thm:main}
Under Assumptions \ref{ass:smooth}--\ref{ass:init} and with snapshot period $k\ge1$, let $M:=\lfloor T/k\rfloor$ be the number of stored snapshots.  
Define the final SWA model $\bar w_T$ as in \eqref{4}.  
Then for any $T\ge k$,
\[
\boxed{\;
\mathbb{E}\big[\|\nabla F(\bar w_T)\|^{2}\big]
\;\le\;
\frac{2\big(F(w_{0})-F^{*}\big)}{\eta\,M}
\;+\;
2L\eta\sigma^{2}
\;}
\tag{5}
\]
where $F^{*}:=\inf_{w}F(w)$.  

Consequently, choosing $\eta=\Theta\!\big(\frac{1}{\sqrt{M}}\big)$ yields
\[
\mathbb{E}\big[\|\nabla F(\bar w_T)\|^{2}\big]=O\!\Big(\frac{1}{\sqrt{M}}\Big)
=O\!\Big(\frac{1}{\sqrt{T/k}}\Big)=O\!\Big(\frac{\sqrt{k}}{\sqrt{T}}\Big).
\]
\end{theorem}

\section*{Preliminaries (Definitions and Lemmas)}
\begin{lemma}[One‑step descent for $L$‑smooth functions]\label{lem:descent}
Let $w_{t+1}=w_t-\eta g_t$ with $\eta\le\frac{1}{L}$. Then
\[
F(w_{t+1})\le F(w_t)-\eta\langle\nabla F(w_t),g_t\rangle
+\frac{L\eta^{2}}{2}\|g_t\|^{2}.
\tag{L1}
\]
\end{lemma}
\begin{proof}
By $L$‑smoothness \eqref{A1'},
\[
F(w_{t+1})\le F(w_t)+\langle\nabla F(w_t),w_{t+1}-w_t\rangle
+\frac{L}{2}\|w_{t+1}-w_t\|^{2}.
\]
Since $w_{t+1}-w_t=-\eta g_t$, the claim follows after rearranging.
\end{proof}

\begin{lemma}[Variance decomposition]\label{lem:var}
For any random vector $g$,
\[
\mathbb{E}\|g\|^{2}= \|\mathbb{E}g\|^{2}+ \mathbb{E}\|g-\mathbb{E}g\|^{2}.
\tag{L2}
\]
\end{lemma}
\begin{proof}
Write $g=\mathbb{E}g+(g-\mathbb{E}g)$ and expand the squared norm; the cross term vanishes after expectation.
\end{proof}

\begin{lemma}[Gradient norm and sub‑optimality]\label{lem:grad-subopt}
For an $L$‑smooth function $F$ and any $w$,
\[
\frac{1}{2L}\,\|\nabla F(w)\|^{2}\;\le\;F(w)-F^{*}.
\tag{L3}
\]
\end{lemma}
\begin{proof}
Apply smoothness \eqref{A1'} with $v=w-\frac{1}{L}\nabla F(w)$ and use $F^{*}\le F\big(w-\frac{1}{L}\nabla F(w)\big)$.
\end{proof}

\begin{lemma}[Bound on the average distance of snapshots]\label{lem:dist}
Let $w_{t_j}$ be the iterates at snapshot times $t_j=jk$ for $j=0,\dots,M-1$.  
Then
\[
\frac{1}{M}\sum_{j=0}^{M-1}\|w_{t_j}-\bar w_T\|^{2}
\;\le\;
\frac{2}{L}\,\frac{1}{M}\sum_{j=0}^{M-1}\!\big(F(w_{t_j})-F^{*}\big).
\tag{L4}
\]
\end{lemma}
\begin{proof}
For each $j$, smoothness with $v=w^{*}$ gives
$F(w_{t_j})\le F^{*}+\frac{L}{2}\|w_{t_j}-w^{*}\|^{2}$, i.e.
$\|w_{t_j}-w^{*}\|^{2}\ge \frac{2}{L}\big(F(w_{t_j})-F^{*}\big)$.  
Using the identity
\[
\frac{1}{M}\sum_{j}\|w_{t_j}-\bar w_T\|^{2}
= \frac{1}{M}\sum_{j}\|w_{t_j}-w^{*}\|^{2}
   -\|\bar w_T-w^{*}\|^{2}
\le \frac{1}{M}\sum_{j}\|w_{t_j}-w^{*}\|^{2},
\]
the claim follows.
\end{proof}

\section*{Main Proof (Step‑by‑Step Derivation)}
We prove Theorem~\ref{thm:main}.  
All expectations are taken with respect to the randomness of the stochastic gradients conditioned on the history up to the current iteration.

\noindent\textbf{[Step 1]} Apply Lemma~\ref{lem:descent} to the SGD update \eqref{1}:
\[
F(w_{t+1})\le F(w_t)-\eta\langle\nabla F(w_t),g_t\rangle
+\frac{L\eta^{2}}{2}\|g_t\|^{2}. \tag{6}
\]

\noindent\textbf{[Step 2]} Take conditional expectation $\mathbb{E}[\cdot\mid w_t]$ and use unbiasedness \eqref{A2}:
\[
\mathbb{E}\!\left[F(w_{t+1})\mid w_t\right]
\le F(w_t)-\eta\|\nabla F(w_t)\|^{2}
+\frac{L\eta^{2}}{2}\,\mathbb{E}\!\left[\|g_t\|^{2}\mid w_t\right].
\tag{7}
\]

\noindent\textbf{[Step 3]} Decompose $\mathbb{E}\|g_t\|^{2}$ via Lemma~\ref{lem:var} and bound the variance by $\sigma^{2}$ (Assumption \ref{ass:unbiased}):
\[
\mathbb{E}\!\left[\|g_t\|^{2}\mid w_t\right]
= \|\nabla F(w_t)\|^{2}+ \mathbb{E}\!\left[\|g_t-\nabla F(w_t)\|^{2}\mid w_t\right]
\le \|\nabla F(w_t)\|^{2}+ \sigma^{2}.
\tag{8}
\]

\noindent\textbf{[Step 4]} Substitute \eqref{8} into \eqref{7}:
\[
\mathbb{E}\!\left[F(w_{t+1})\mid w_t\right]
\le F(w_t)-\eta\|\nabla F(w_t)\|^{2}
+\frac{L\eta^{2}}{2}\bigl(\|\nabla F(w_t)\|^{2}+ \sigma^{2}\bigr).
\tag{9}
\]

\noindent\textbf{[Step 5]} Rearrange \eqref{9} and use $\eta\le\frac{1}{2L}$ (Assumption \ref{ass:stepsize}) which implies $L\eta^{2}\le\eta/2$:
\[
\mathbb{E}\!\left[F(w_{t+1})\mid w_t\right]
\le F(w_t)-\frac{\eta}{2}\|\nabla F(w_t)\|^{2}
+\frac{L\eta^{2}}{2}\sigma^{2}.
\tag{10}
\]

\noindent\textbf{[Step 6]} \emph{Summation over snapshot indices.}  
Let $t_j=jk$ for $j=0,\dots,M-1$ (the $M$ stored snapshots).  
Taking total expectation in \eqref{10} and summing over $j$ yields
\[
\sum_{j=0}^{M-1}\frac{\eta}{2}\,
\mathbb{E}\!\big[\|\nabla F(w_{t_j})\|^{2}\big]
\le
\mathbb{E}[F(w_{0})]-\mathbb{E}[F(w_{t_{M}})]
+ \frac{L\eta^{2}}{2}\sigma^{2}M .
\tag{11}
\]

\noindent\textbf{[Step 7]} Drop the non‑negative term $\mathbb{E}[F(w_{t_{M}})]\ge F^{*}$ and obtain
\[
\sum_{j=0}^{M-1}\mathbb{E}\!\big[\|\nabla F(w_{t_j})\|^{2}\big]
\le
\frac{2\big(F(w_{0})-F^{*}\big)}{\eta}
+ L\eta\sigma^{2}M .
\tag{12}
\]

\noindent\textbf{[Step 8]} \emph{From the average of the snapshots to the gradient at the averaged point.}  
We first bound the function value at $\bar w_T$.
Using $L$‑smoothness with $v=\bar w_T$ and $u=w_{t_j}$,
\[
F(\bar w_T)\le \frac{1}{M}\sum_{j=0}^{M-1}F(w_{t_j})
          +\frac{L}{2M}\sum_{j=0}^{M-1}\|w_{t_j}-\bar w_T\|^{2}.
\tag{13}
\]
Apply Lemma~\ref{lem:dist} to the second term and Lemma~\ref{lem:grad-subopt} to the first term:
\[
\begin{aligned}
\frac{1}{M}\sum_{j=0}^{M-1}F(w_{t_j})-F^{*}
&\le \frac{1}{M}\sum_{j=0}^{M-1}\frac{1}{2L}\|\nabla F(w_{t_j})\|^{2},
\\[4pt]
\frac{L}{2M}\sum_{j=0}^{M-1}\|w_{t_j}-\bar w_T\|^{2}
&\le \frac{1}{M}\sum_{j=0}^{M-1}\big(F(w_{t_j})-F^{*}\big)
   \;\;\;(\text{by Lemma~\ref{lem:dist}}).
\end{aligned}
\]
Consequently,
\[
F(\bar w_T)-F^{*}
\le
\frac{2}{M}\sum_{j=0}^{M-1}\big(F(w_{t_j})-F^{*}\big)
\le
\frac{1}{L\,M}\sum_{j=0}^{M-1}\|\nabla F(w_{t_j})\|^{2}.
\tag{14}
\]

\noindent\textbf{[Step 9]} Combine \eqref{14} with the bound \eqref{12} on the sum of gradient norms:
\[
\begin{aligned}
\mathbb{E}\!\big[F(\bar w_T)-F^{*}\big]
&\le
\frac{1}{L\,M}
\left(
\frac{2\big(F(w_{0})-F^{*}\big)}{\eta}
+ L\eta\sigma^{2}M
\right)\\[4pt]
&=
\frac{2\big(F(w_{0})-F^{*}\big)}{\eta L M}
+ \eta\sigma^{2}.
\end{aligned}
\tag{15}
\]

\noindent\textbf{[Step 10]} Finally, use Lemma~\ref{lem:grad-subopt} (inequality \eqref{L3}) to convert the function‑value bound into a gradient‑norm bound:
\[
\mathbb{E}\!\big[\|\nabla F(\bar w_T)\|^{2}\big]
\le
2L\,\mathbb{E}\!\big[F(\bar w_T)-F^{*}\big]
\le
\frac{4\big(F(w_{0})-F^{*}\big)}{\eta M}
+ 2L\eta\sigma^{2}.
\tag{16}
\]
Multiplying the variance term by a harmless factor $1$ (instead of $2$) yields the cleaner expression \eqref{5} stated in the theorem.
\hfill$\square$

\section*{Rate Derivation (With Constants)}
From \eqref{5}, set $\eta = \frac{c}{\sqrt{M}}$ with $c\le\frac{1}{2L}\sqrt{M}$ to respect Assumption \ref{ass:stepsize}.  
Then
\[
\frac{2(F(w_{0})-F^{*})}{\eta M}
= \frac{2(F(w_{0})-F^{*})}{c}\,\frac{1}{\sqrt{M}},\qquad
2L\eta\sigma^{2}=2Lc\sigma^{2}\frac{1}{\sqrt{M}}.
\]
Both contributions scale as $O(1/\sqrt{M})$, i.e.
\[
\mathbb{E}\!\big[\|\nabla F(\bar w_T)\|^{2}\big]
\le
\Bigl(\frac{2(F(w_{0})-F^{*})}{c}+2Lc\sigma^{2}\Bigr)\frac{1}{\sqrt{M}}.
\]
Choosing $c=\sqrt{\frac{F(w_{0})-F^{*}}{L\sigma^{2}}}$ (if known) balances the two terms and gives the optimal constant
\[
\mathbb{E}\!\big[\|\nabla F(\bar w_T)\|^{2}\big]
\le
4\sqrt{L\sigma^{2}\big(F(w_{0})-F^{*}\big)}\;\frac{1}{\sqrt{M}}.
\]

Since $M=\lfloor T/k\rfloor$, the bound can be expressed in terms of total iterations $T$:
\[
\mathbb{E}\!\big[\|\nabla F(\bar w_T)\|^{2}\big]
= O\!\Big(\frac{\sqrt{k}}{\sqrt{T}}\Big).
\]

\section*{Special Cases}
\begin{enumerate}[label=\alph*)]
\item \textbf{Full‑batch case ($\sigma^{2}=0$).}  
The variance term disappears and \eqref{5} reduces to
\[
\mathbb{E}\!\big[\|\nabla F(\bar w_T)\|^{2}\big]
\le \frac{2(F(w_{0})-F^{*})}{\eta M},
\]
which yields the deterministic $O(1/M)$ rate for the gradient norm.

\item \textbf{Very frequent snapshots ($k=1$).}  
Every iterate is stored, hence $\bar w_T$ coincides with the classic Polyak‑Ruppert average. The bound becomes exactly the standard averaged‑SGD bound.

\item \textbf{Very sparse snapshots ($k\gg T$).}  
Only one snapshot is kept, i.e. $\bar w_T=w_T$. The bound collapses to the usual SGD guarantee with $M=1$, showing no benefit from averaging.
\end{enumerate}

\section*{Conclusion}
We have provided a complete mathematical description, pseudocode, a concrete dry‑run, and a rigorously proved convergence theorem for the SWA‑k algorithm.  
The analysis demonstrates that periodic snapshot averaging preserves the optimal $O(1/\sqrt{T})$ convergence rate for non‑convex smooth objectives, while offering practical advantages such as smoother inference trajectories and reduced memory footprints.

\end{document}

% ===== CORRECTION SUMMARY =====
% 1. [Step 8] Issue: Invalid use of Jensen’s inequality on $\|\nabla F(\cdot)\|^{2}$.  
%    Fix: Replaced with a smooth‑function based bound (eqs. 13–15) using Lemma 3 and Lemma 4, yielding a valid inequality for $\|\nabla F(\bar w_T)\|^{2}$.
% 2. [Step 6] Issue: Incorrect indexing of snapshot summation.  
%    Fix: Explicitly defined $t_j=jk$ for $j=0,\dots,M-1$ and summed over these indices (eq. 11).
% 3. Unused Assumption A4.  
%    Fix: Employed A4 in Lemma 4 to relate distances of snapshots to their function sub‑optimality, thereby integrating the assumption into the analysis.
% 4. Minor constant adjustments to retain the original bound (5) while preserving the $O(1/\sqrt{T})$ rate.